\section{Persistence}
\label{sec:Persistence}

The Archimesh server persists all model state to Neo4j. This chapter describes
the graph model, repository layer and key invariants.

\subsection{Repository layer}

The persistence layer is organized into a split repository with a main
coordinator and specialized support classes:

\begin{description}
\item[\code{Neo4jRepositoryImpl}] Main coordinating repository. Orchestrates
  op-log writes, materialized state mutations, head revision updates,
  snapshots, rebuilds, imports, exports, compactions and integrity checks.
\item[\code{Neo4jOpLogSupport}] Commit and operation persistence. Writes
  \code{:Commit} and \code{:Op} nodes, maintains the \code{NEXT}
  chain.
\item[\code{Neo4jMaterializedStateSupport}] Materialized graph mutation.
  Creates, updates and deletes elements, relationships, views, view objects,
  connections and folders. Applies LWW merges and tombstone checks.
  Maintains property clocks.
\item[\code{Neo4jReadSupport}] Snapshot export, checkout delta reads,
  consistency checks. Reads materialized state and converts it to the JSON
  snapshot/export format.
\item[\code{Neo4jCompactionSupport}] Bounded pruning of published
  operation payloads covered by a verified checkpoint and retention watermark.
  Retains acceptance receipts, clocks, and tombstones.
\item[\code{Neo4jCheckpointSupport}] Lossless raw graph capture, schema
  and checksum verification, and checkpoint restoration.
\end{description}

\subsection{Graph model}

The database stores these kinds of state per model:

\begin{enumerate}
\item Model catalog and ACL metadata
\item Materialized current state
\item Committed op-log history
\item Merge metadata (tombstones, property clocks)
\item Immutable recovery checkpoints and named tag snapshots
\item Durable publication intents and retained acceptance receipts
\end{enumerate}

All identity is model-scoped: materialized entity identity is
\code{(modelId, id)}; commit identity is \code{(modelId, opBatchId)};
tag identity is \code{(modelId, tagName)}.

\subsubsection{Core model node}

Every model is rooted at a \code{:Model} node:

\begin{lstlisting}[caption={Model node properties}]
{
  modelId, modelName,
  registered: true,
  headRevision: <number>,
  createdAt, updatedAt,
  adminUsers: ["user1", ...],
  writerUsers: ["user2", ...],
  readerUsers: ["user3", ...]
}
\end{lstlisting}

ACL data is stored directly on \code{:Model} as user-id arrays.

\subsubsection{Materialized state nodes}

\begin{longtable}{p{30mm}p{124mm}}
\toprule
Label & Properties \\
\midrule
\code{:Element} & \code{modelId}, \code{id}, \code{archimateType},
  \code{name}, \code{documentation},
  \code{name\_lamport}, \code{name\_clientId},
  \code{documentation\_lamport}, \code{documentation\_clientId} \\
\code{:Relationship} & \code{modelId}, \code{id}, \code{archimateType},
  \code{name}, \code{documentation}, \code{sourceId}, \code{targetId},
  plus LWW clocks for name, documentation, source and target. \\
\code{:Folder} & \code{modelId}, \code{id}, \code{folderType},
  \code{name}, \code{documentation}, plus LWW clocks. \\
\code{:View} & \code{modelId}, \code{id}, \code{name},
  \code{documentation}, \code{notationJson}, plus LWW clocks. \\
\code{:ViewObject} & \code{modelId}, \code{id}, \code{notationJson},
  plus per-field LWW clocks for geometry, style and label fields. \\
\code{:Connection} & \code{modelId}, \code{id},
  \code{sourceViewObjectId}, \code{targetViewObjectId},
  \code{representsId}, \code{notationJson}, plus per-field LWW clocks. \\
\bottomrule
\end{longtable}

\subsubsection{Materialized relationships}

\begin{lstlisting}[caption={Core relationships}]
(:Model)-[:HAS_ELEMENT]->(:Element)
(:Model)-[:HAS_REL]->(:Relationship)
(:Relationship)-[:SOURCE]->(:Element)
(:Relationship)-[:TARGET]->(:Element)

(:Model)-[:HAS_FOLDER]->(:Folder)
(:Folder)-[:HAS_FOLDER]->(:Folder)
(:Folder)-[:CONTAINS_ELEMENT]->(:Element)
(:Folder)-[:CONTAINS_REL]->(:Relationship)
(:Folder)-[:CONTAINS_VIEW]->(:View)

(:Model)-[:HAS_VIEW]->(:View)
(:View)-[:CONTAINS]->(:ViewObject)
(:ViewObject)-[:REPRESENTS]->(:Element)
(:ViewObject)-[:CHILD_MEMBER]->(:ViewObject)

(:View)-[:CONTAINS]->(:Connection)
(:Connection)-[:REPRESENTS]->(:Relationship)
(:Connection)-[:FROM]->(:ViewObject)
(:Connection)-[:TO]->(:ViewObject)
\end{lstlisting}

\subsubsection{Folder rules}

\begin{itemize}
\item Root folders use stable synthetic ids:
  \code{folder:root-strategy}, \code{folder:root-business},
  \code{folder:root-diagrams}, etc.
\item User folders use immutable ids --- rename changes the name property
  but not the node identity.
\item A model-tree object belongs to exactly one folder at a time.
\item Cross-root folder moves are rejected by validation.
\item Non-empty folder delete is rejected.
\end{itemize}

\subsubsection{Op-log history}

Commit nodes preserve the accepted timeline. Linked operation payloads remain
until checkpoint-covered pruning:

\begin{lstlisting}[caption={Op-log nodes}]
(:Commit {
  modelId, opBatchId,
  revisionFrom, revisionTo,
  ts, opCount,
  published, payloadPruned, batchJson
})

(:Op {
  seq, type, targetId, payloadJson
})
\end{lstlisting}

Relationships:

\begin{lstlisting}[caption={Op-log relationships}]
(:Model)-[:HAS_COMMIT]->(:Commit)
(:Commit)-[:HAS_OP]->(:Op)
(:Commit)-[:NEXT]->(:Commit)
\end{lstlisting}

\begin{itemize}
\item \code{revisionFrom..revisionTo} is the assigned revision range for
  one committed submit batch.
\item \code{payloadJson} stores the normalized op payload.
\item \code{seq} preserves op order inside a batch.
\item \code{NEXT} keeps the timeline linear.
\item Production atomic commits store \code{batchJson} and
  \code{published=false} together with projection and head changes.
  Broker acknowledgment permits setting \code{published=true}.
\item Pruning removes \code{Op} nodes and \code{batchJson}, then marks
  \code{payloadPruned=true}. The commit remains as an acceptance receipt.
\end{itemize}

\subsubsection{Recovery checkpoints}

\begin{lstlisting}[caption={Checkpoint node and ownership}]
(:Model)-[:HAS_CHECKPOINT]->(:ModelCheckpoint {
  modelId, revision, sha256, payloadJson
})
\end{lstlisting}

\code{payloadJson} contains the versioned checkpoint document, including
creation timestamp, raw nodes and edges, all merge metadata and the checksum.
Uniqueness is \code{(modelId, revision)}. At an existing boundary, creation
returns the immutable checkpoint rather than replacing it.

The checkpoint's projection does not duplicate catalog metadata, ACLs, tags,
commits or its own storage node. Its reserved model endpoint is rebound to the
destination model on restore. A checkpoint is richer than a client snapshot;
it includes raw clock and orphan safety state needed for deterministic merges.
See Section~\ref{sec:Checkpoints} for validation and recovery requirements.

\subsubsection{Tags}

Tags are stored as standalone nodes:

\begin{lstlisting}[caption={Tag node}]
(:ModelTag {
  modelId, tagName, description,
  revision, createdAt, snapshotJson
})
\end{lstlisting}

\begin{itemize}
\item Tags are immutable --- once created, they cannot be modified or deleted
  (tag deletion is disabled by default).
\item \code{snapshotJson} stores the materialized snapshot at the tagged
  revision.
\item Tags do not create branches; they are read-only pointers on the linear
  timeline.
\end{itemize}

\subsubsection{Merge metadata --- tombstones}

Tombstones prevent stale recreates after deletes:

\begin{lstlisting}[caption={Tombstone nodes}]
(:ElementTombstone { modelId, id, lamport, clientId })
(:RelationshipTombstone { modelId, id, lamport, clientId })
(:ViewTombstone { modelId, id, lamport, clientId })
(:ViewObjectTombstone { modelId, id, lamport, clientId })
(:ConnectionTombstone { modelId, id, lamport, clientId })
\end{lstlisting}

Tombstones are anchored from \code{:Model} via specific relationship types
(\code{HAS\_ELEMENT\_TOMBSTONE}, etc.). Compaction reports tombstone
eligibility but retains them for delete safety.

\subsubsection{Merge metadata --- property clocks}

Property merge state is stored in dedicated metadata nodes:

\begin{lstlisting}[caption={PropertyClock node}]
(:PropertyClock {
  modelId, targetId, key,
  lamport, clientId, deleted
})
\end{lstlisting}

Anchored via \code{(:Model)-[:HAS\_PROPERTY\_CLOCK]->(:PropertyClock)}.
Supports LWW and OR-set style property semantics. Compaction retains all
property clocks, including orphan safety metadata.

\subsection{Neo4j constraints}

The schema file at \filepath{neo4j/schema.cypher} defines the following
uniqueness constraints:

\begin{lstlisting}[caption={Uniqueness constraints}]
(:Model {modelId})                    UNIQUE
(:Element {modelId, id})              UNIQUE
(:Relationship {modelId, id})         UNIQUE
(:View {modelId, id})                 UNIQUE
(:Folder {modelId, id})               UNIQUE
(:ViewObject {modelId, id})           UNIQUE
(:Connection {modelId, id})           UNIQUE
(:Commit {modelId, opBatchId})        UNIQUE
(:ModelTag {modelId, tagName})        UNIQUE
(:ModelCheckpoint {modelId, revision}) UNIQUE
\end{lstlisting}

An index on \code{(:Commit)} by \code{(modelId, revisionFrom, revisionTo)}
supports efficient delta window queries.

\subsection{Persistence invariants}

\begin{itemize}
\item Materialized entity identity is \code{(modelId, id)} --- never rely
  on node-internal ids across queries.
\item Commit/idempotency identity is \code{(modelId, opBatchId)} ---
  duplicate submission of the same batch is a no-op.
\item Atomic model commits persist log, projection, head and delivery intent
  together. A failed transaction changes none of them.
\item Recovery requires a verified checkpoint and complete tail, or full log
  from revision 1. Replacement is transactional. Pruning keeps acceptance
  receipts and all CRDT safety metadata.
\item Repository write paths must fail fast on Neo4j persistence errors ---
  never silently skip a mutation.
\item Notation field definitions must flow through
  \code{NotationMetadata}, the shared source of truth for supported
  notation fields and persisted merge metadata.
\item Compaction never allows stale \code{Create*} or \code{Update*} to
  beat a logically newer delete --- tombstones are retained.
\end{itemize}

\subsection{Configuration}

\begin{longtable}{p{42mm}p{50mm}p{62mm}}
\toprule
Property & Env var & Default \\
\midrule
\code{app.neo4j.uri} & \code{NEO4J\_URI} &
  \code{bolt://localhost:7687} \\
\code{app.neo4j.username} & \code{NEO4J\_USER} & \code{neo4j} \\
\code{app.neo4j.password} & \code{NEO4J\_PASSWORD} &
  \code{devpassword} \\
\bottomrule
\end{longtable}

In production, always set \code{NEO4J\_PASSWORD} to a strong value. The
default \code{devpassword} is for local development only.
